\section{Preliminary Experimental Analysis}
\label{sec:evaluations}

%The proposed LDPC decoder was simulated using MATLAB’s nonstiff, single step, 5th-order differential solver (ode45). To obtain solution quality for LDPC decoding using the QUBO formulation, we used the original BRIM model described in paper~\cite{afoakwa.hpca21} using Matlab.

Both formulations QUBO and modified QUBO were tested using Simulated annealing algorithm ~\cite{Kirkpatrick671} for solution quality. LDPC codes generated by 5G-NR standards are used as benchmarks. 

\subsection{Comparison of solution quality with QUBO Formulation }\label{sec:solution_quality}

One thousand messages were encoded using different LDPC codes to obtain the solution quality. These codes were constructed using the 5G-NR standard's base graph 1 with varying factors of expansion. The encoded messages were then modulated using BPSK modulation. The noise was added to the modulated signals based on the signal-to-noise ratio (SNR). These corrupted signals were used to test the different algorithms for solution quality. 

Figure~\ref{fig:ber} shown the solution quality of LDPC decoding using QUBO formulation and proposed formulation (modified QUBO). Both modified QUBO and QUBO were annealed using simulated annealing algorithm for 1000 iterations to obtain BER. As the simulated annealing algorithm is  probabilistic we used a different metric called Expected BER to estimate the solution quality. Each message (or instance) was annealed 100 times with random initializations. Each different anneal gives us a different solution. We find the distinct solutions in these different anneals and rank them according to their Ising energies. Considering this order static, the \textit{expected} Bit Error rate (E(BER($N_a$)))~\cite{kim2019leveraging} for each instance is calculated using the formula~\ref{eqn:expected_ber}.

\begin{equation}\label{eqn:expected_ber}
\begin{split}
E(BER(N_a)) = \sum_{q=1}^{L} \left[ (\sum_{r=q}^{L} P(r))^{N_a} -  (\sum_{r=q+1}^{L} P(r))^{N_a}\right] \\ F(q)/K    
\end{split}
\end{equation}

where K is the number of bits in the message, L is the number of unique solutions obtained by different anneals, r is the rank index of each solution, $P(r)$ is the probability of obtaining the $r_{th}$ solution, $F$ is the number of bit errors in the $q^{th}$ solution, relative to the ground truth and $N_a$ represents number of anneals. The average of 1000 different expected BER is reported as the final Bit Error rate for each SNR. As shown in the figure, modified QUBO gives us the better BER when compared to QUBO formulation.

\begin{figure}[htb]
  \centering
 \begin{minipage}[b]{.40\textwidth}
   \includegraphics[width =\textwidth]{LDPC codes/FIGS/expf_16.pdf}
  %\caption{This is first figure}
  \end{minipage}%
  \hfill
  \begin{minipage}[b]{.5\textwidth}
  \includegraphics[width = 0.4\textwidth]{LDPC codes/FIGS/expf_2.pdf} 
  \includegraphics[width = 0.4\textwidth]{LDPC codes/FIGS/expf_4.pdf}
  %\caption{caption}
  \end{minipage} 
  \caption{Bit Error rate achieved by different algorithms at different Signal-noise ratios. Lower the better.}
\label{fig:ber}
\end{figure}

\section{Things need to be done}

\subsection{Determine accurate simulation of hardware}

The proposed hardware can be modeled using cadence but it takes lot of time to simulate. Hence we want to faithfully model the hardware using differential equation and solve it using MATLAB.

\subsection{Determine the architecture of the system}
There are two ways to determine the architecture \ding{172} to design the system to accommodate all LDPC codes determine by the standard 
\ding{173} to design the system for a particular parity check matrix
we want to explore how the system performance varies accordingly.

\subsection{Evaluate our system with state of the art LDPC decoders}

There exists plenty of literature about designing LDPC decoders. Firstly, we want to investigate if there is any advantage in decoding LDPC in our method compared to the Min-Sum algorithm. Secondly, the existing decoders are designed for different purposes using different technologies. We want to fairly compare all these decoders by adapting their design to solve a common problem and investigate how our design would compare against them.

\subsection{Noise analysis}

As our design is analog, we are prone to noise. Hence we want to investigate how our system's performance would get affected by noise. 

%\todo{To be implemented}
%\subsection{Circuit Implementation using Cadence}
%\todo{To be implemented}

